\documentclass[../main]{subfiles}
\begin{document}
\section{Corte y Sección}

\imgfull{images/09/corte1}{Corte.}

 El dibujo técnico debe dar una idea clara y precisa del exterior del objeto representado y también de sus características interiores. Mientras menos líneas ocultas (trazos cortos) posea el dibujo, será más fácil de interpretar; esto se logra aplicando cortes y secciones.

\img{images/09/corte4}{Señalización de corte.}

Se representa pasando un plano imaginario a la pieza, la dividirá en dos partes y representamos lo que está en el plano cortante (parte que se representará rayada) y lo que está detrás del plano cortante (parte que se representará no rayada).


\subsubsection{Planos de cortes}

 \imgfull{images/09/corte3}{Ejemplo de corte.}

Los planos de corte se representan mediante líneas de trazo largos y trazos cortos, cuyos extremos se dibujarán con trazos gruesos. Las líneas de corte llevarán en sus extremos flechas que se anteponen a la línea de corte indicando la dirección y sentido visual.

La denominación del corte se indica con letra mayúscula en los extremos de la línea de corte, en la posición de lectura normal, sobre la línea de la flecha o en el costado de ella, como se ve en todas las figuras.

\subsubsection{Corte quebrado}
La línea de indicación del corte podrá ser recta, quebrada o curva En caso de plano de corte quebrado, en cada quiebre se colocará una letra mayúscula siguiendo el orden alfabético, siendo la letra final igual a la inicial, en el trazo en el áncgulo se engrosará El corte quebrado se emplea cuando hay varios detalles internos no alineados, quebrando en el plano de corte. Deste modo se efectúa una sola vista de corte.

\imgfull{images/09/corte6}{Corte quebrado.}

\subsubsection{Rayado del corte}

El rayado de la sección se realiza con líneas finas igualmente separadas con inclinación de $45^o$ respecto al eje principal como se puede ver en la figura.

\imgfull{images/09/corte5}{Trazado de corte.}


\newpage
\subsubsection{Actividades}
\actividad{Leer, responder y dibujar}
\begin{enumerate}
    \item ¿Qué es corte?
    \item ¿En un corte quebrado como puede ser la indicación del corte?
    \item ¿Cómo se realiza el rayado de un corte?
    \item Copiar en la siguiente hoja normalizada completar el rótulo y acotar.
    \imgfull{images/09/corte2}{Copiar en la siguiente página.}
\end{enumerate}
\end{document}